\DocumentMetadata{pdfversion=2.0}
\documentclass[conference]{IEEEtran}
% \IEEEoverridecommandlockouts
% The preceding line is only needed to identify funding in the first footnote. If that is unneeded, please comment it out.
\usepackage{cite}
\usepackage{amsmath,amssymb,amsfonts}
\usepackage{algorithmic}
\usepackage{graphicx}
\usepackage{textcomp}
\usepackage{xcolor}
\usepackage[colorlinks=true,urlcolor=blue,linkcolor=blue]{hyperref} % To allow Alt Text and Linking for Headers, and to allow blue, underlined URL hyperlinks

\def\BibTeX{{\rm B\kern-.05em{\sc i\kern-.025em b}\kern-.08em
    T\kern-.1667em\lower.7ex\hbox{E}\kern-.125emX}} % Previous lines from IEEE template
\begin{document}

\title{Semantic Analysis of Domain Data for Proactive Identification of Cybersecurity Indicators}

\author{\IEEEauthorblockN{Nic Recasens}
\IEEEauthorblockA{\textit{Computer Science} \\
\textit{Georgia Southern University}\\
Statesboro, Georgia, USA \\
jr38088@georgiasouthern.edu}
\and
\IEEEauthorblockN{Nigel Smith}
\IEEEauthorblockA{\textit{Computer Science} \\
\textit{Georgia Southern University}\\
Statesboro, Georgia, USA \\
ns15468@georgiasouthern.edu}}

\maketitle

%\section{Abstract}
% 150 - 250 words

\section{Introduction and Problem Statement}
%@j-recasens @ns15468-gasou research question has to be added explicitly below for full points
In the modern world, attackers are a step ahead of all Cybersecurity professionals. 
Professionals need to be able to adapt to every threat out there, while attackers only need to be able to adapt to a single fatal weakness.
As such, the realm of machine learning and predictive analytics is of known importance in Cybersecurity research.
For decades, machine learning analysis has helped keep defenders on top of their game by preventing novel attacks from gaining footholds.
In the 90s, when the number of attackers was smaller, professionals were able to get away with only blocking known threats.
But with Ransomware and Malware becoming billion dollar multi-national black markets, we need to be able to apply risk trends to novel data.
At the forefront of this prevention is the identification and blocking of spam and phishing.
An individual datum related to spam and phishing is called an indicator of attack (IoA), and they let us know if someone is targeted by an attacker.
And if they are used successfully, they become indicators of compromise (IoC).
IoCs have the negative reputation of being indicative of failed Cybersecurity protection, but in relevance to this project, they are the second half of the puzzle that is semantic analysis of indicators.
Via the correlation of IoAs and/or IoCs, numerous researchers have pushed the boundary forward over the decades since the aforementioned 1990s.
The attackers may be numerous in their successes, but they leave behind their own value as well: data. 
With that data, we can see how pervasive cyber threats are, train models to simulate said attackers, and allow the very computers they attack help mitigate the risks themselves. 
Through this literature review, we will go over 8 papers that show how researchers achieved these goals, and show through their research how important cyber-risk mitigation truly is in our modern interconnected world.

\section{Related Literature}

\subsection{Literature Selection Methodology}
%% Explain the process of selecting articles, detailing:
%Keywords used
At the most basic level, the 3 keywords we used to guide our literature selection process were "spam", "phishing", and "semantic".
Spam and Phishing are the strongest indicators Cybersecurity professionals have as indicators of attack.
In fact, Business Email Compromise or (BEC) can be identified as the largest threat to the modern digital business \cite{saidani1}.
As such, once we have papers scoped to the huge risk spam and phishing compose, the "semantic" keyword allows us to predominantly find Computer Science papers detailing how semantic analysis can help with this cause.
Some linguistic journals do appear with these keywords, but with the current rise of the Large-Language Model, semantic research is seeing a heavy Computer Science edge.

%Databases accessed
For searching these terms, Google Scholar was the primary resource, though GALILEO also was utilized for accessing some non-free files. When on our Georgia Southern University network, some research journals allow our IP to access full papers without logon, which was also useful for reading full texts. Publication Year and citation information was considered and prioritized past the above keyword filtering, and was based on the Google Scholar data for those indicators.

%and Inclusion and exclusion criteria
And to wrap up our selection process, our primary exclusion criteria was mentioned earlier, and that is a Computer Science focus.
Some articles were industry paid, so were more like Sociology in nature.
And for other articles, they were focusing on the behavioral aspects of phishing and spam, which while useful, are not in scope for this review.
And speaking of scope, the primary \textit{inclusion} criteria was related to this course's guidelines, and that is publication year.
Many useful articles at the beginning of semantic analysis for spam prevention were made in the late 2000's, but they were specifically asked to be ignored for this paper by the paper requirements.
As such, only 'recent' research was provided, meaning publication dates past 2020 as much as it could be helped.
These criteria, alongside with the above requirements and toolset mindsets, allowed us to choose the below 8 papers to systematically review.

\subsection{Reactive Detection of Attacks}
The earliest theme addresses detection of spam and phishing as it occurs. Saidani, Adi, and Allili \cite{saidani1} confront the fact that spam continues to account for over 56\% of total email traffic, yet traditional detection systems largely rely on the provenly ineffectual "bag-of-words" model, which evaluates emails as an unstructured collection of independent words without considering grammar, context, or word order. Spammers purposefully exploit the weaknesses of basic keyword filters through obfuscation, motivating their proposed two-level semantic classification framework, which achieved up to 98\% accuracy across domain-specific classifiers \cite{saidani1}. Sonowal and Gunikhan \cite{sonowal1} extend this challenge to SMS, noting that unlike email's 20-30\% open rates, SMS messages boast a staggering 98\% open rate, making smishing a disproportionately dangerous vector. Their feature-based approach reduced 52 input dimensions to just 20 via Kendall rank correlation --- a 61.53\% decrease --- while maintaining 98\% accuracy, a reduction critical for viability on resource-constrained mobile devices \cite{sonowal1}. Both papers share a common vulnerability: as generative AI tools become more indistinguishable when mimicking professional readability, the semantic dissonance these systems rely on may diminish \cite{sonowal1}.

\subsection{Proactive Data Gathering}
A second thread in the literature pushes beyond reactive detection toward proactively building threat intelligence. Zhang et al. \cite{zhang1} identify that before spam and phishing can be remediated, we have to know what they actually \textit{are} via data collection, noting that conventional CTI relies on "fixed-point" monitoring of curated feeds that is time-consuming, delayed, and incomplete. Their iMCircle system addresses this with an active recursive architecture that mines IoCs continuously from open-source web data using NLP and SVM classification, achieving 96.70\% precision and an F1 score of 90.46\% \cite{zhang1}. However, the system remains entirely dependent on publicly indexed information and cannot surface threats not yet discussed on the open web \cite{zhang1}. Zieni, Massari, and Calzarossa \cite{zieni_phishing_2023} complement this by surveying the broader detection landscape, finding not a single technical gap but a fragmentation of approaches across list-based, similarity-based, and ML-based methods. Their survey identifies Random Forest as the most consistently performant traditional algorithm, while also flagging a critical research gap: URL shortening services render most lexical URL features meaningless \cite{zieni_phishing_2023}.

\subsection{Proactive Prevention of Attacks}
The most recent papers in the review target prevention rather than detection, driven by capabilities unlocked by large language models and self-updating ML systems. Brezeanu, Archip, and Artene \cite{brezeanu_phish_2025} observe that phishing kits --- reusable HTML templates allowing attackers to rapidly deploy fraudulent pages --- frequently reuse identical HTML structures across campaigns, making the underlying code structure a more reliable detection signal than visual or textual content. Their Phish Fighter system applies agglomerative hierarchical clustering and a self-retraining Random Forest classifier, achieving weighted precision, recall, and F1 scores exceeding 90\% \cite{brezeanu_phish_2025}. Desolda, Greco, and Viganò \cite{desolda_apollo_2025} address a parallel problem: that no existing system achieves perfect accuracy, meaning borderline cases reach the user, and the warnings users actually see in browsers and email clients are widely ineffective. Their APOLLO system, powered by GPT-4o, achieved 97.4\% accuracy without hand-crafted features, rising to 99.9\% when mature VirusTotal data was available \cite{desolda_apollo_2025}. Critically, both systems still encounter the zero-day ceiling: Phish Fighter requires a minimum of three pages from a new kit before clustering, and APOLLO's performance worsens when VirusTotal returns uncertain data on newly minted URLs \cite{brezeanu_phish_2025, desolda_apollo_2025}.

\subsection{Psychological Studies}
No technical system operates in a vacuum. Dash and Ansari \cite{dash1} note that despite advancements in defensive systems, "naïve information security practices" account for between 72\% and 95\% of all cyber threats, motivating their study of AI-enabled training tools like viCyber, which can provide real-time feedback and adapt to each user's understanding level \cite{dash1}. Sarno, Harris, and Black \cite{dawn_m_sarno_which_2023} ground this in empirical data, finding that some phishing emails will inevitably reach the inbox, at which point the user becomes the last line of defense. Their regression study of 150 participants found that younger adults were most susceptible --- directly contradicting the popular assumption that older adults are the most vulnerable population --- and that self-reported phishing experience may foster overconfidence rather than genuine detection skill \cite{dawn_m_sarno_which_2023}. Together, these papers reinforce that interpretable outputs and end-user empowerment are not secondary concerns but core requirements for any detection system.

\subsection{Overall Literature Synthesis and Conclusion}
Two common themes reveal themselves across these papers. First, existing tools are commonly vulnerable to zero-day attacks, which suggest the need for robust ML and deep learning tools. From the reactive blocklists reviewed by Zieni et al.\cite{zieni_phishing_2023} to the phishing kit clustering of Brezeanu et al.\cite{brezeanu_phish_2025} to the retrieval-augmented classification of APOLLO\cite{desolda_apollo_2025}, every system reviewed encounters the same fundamental constraint: detection accuracy degrades sharply when the threat has not been seen before. Brezeanu's Phish Fighter requires a minimum of three pages from a new kit before it can form a cluster, and APOLLO's performance actually worsens when VirusTotal returns uncertain data on newly minted URLs. Even iMCircle, which was designed specifically for continuous discovery, can only surface indicators that have already been publicly discussed somewhere on the open web\cite{zhang1}. The zero-day problem is not merely a limitation noted in passing by these authors; it is the central unsolved challenge that connects the entire body of work reviewed here.

Second, the ability to meet this need exists, but is still in its infancy, with effective ML tools having only existed for a matter of a few years. The progression visible across the papers is striking: Saidani et al.\cite{saidani1} advanced past bag-of-words with domain-specific semantic rules in 2020, Sonowal and Gunikhan\cite{sonowal1} showed that careful feature engineering could make ML viable even on resource-constrained mobile devices, and by 2025, APOLLO demonstrated that a general-purpose LLM could classify phishing emails at 97.4\% accuracy without any hand-crafted features at all\cite{desolda_apollo_2025}. 

Our project would aim to reconcile the findings and shortcomings of these papers by utilizing cutting-edge ML tools for the purpose of highly accurate and immediate phishing detection. Specifically, the research gaps identified here point toward a system that can operate effectively against novel threats without requiring a warm-up period of previously seen examples, that can leverage the semantic reasoning capabilities demonstrated by LLM-based approaches while maintaining the structural analysis strengths shown by methods like Phish Fighter, and that produces outputs interpretable enough to support the human end user who remains, as this review has shown, the last and most variable line of defense.

\section{Data Processing and Framework Diagram}
All data is stored on Google Cloud Storage Containers, primarily to achieve the goal that our results can be replicated by our professor.
% including wrangling, cleaning, and exploratory visualizations
%% preliminary results, along with dataset details are required for revised review

%\section{Methodology}
% including DS/ML techniques used with justifications

%\section{metrics, model performance, visualizations}
% including wrangling, cleaning, and exploratory visualizations

%\section{Discussion}
% including interpretation, limitations, implications

%\section{Conclusion and Potential Future Work}

% @j-recasens As an FYI, I've learned that even if a source is in 
% project.bib, it won't show up in the compiled code without being \cite'd first.
\bibliographystyle{plain} 
\bibliography{project}

\end{document}




%%%% Help Messages from IEEE (moved to the bottom to not be lost):

%% \subsection{Identify the Headings}

% Headings, or heads, are organizational devices that guide the reader through 
% your paper. There are two types: component heads and text heads.

% Component heads identify the different components of your paper and are not 
% topically subordinate to each other. Examples include Acknowledgments and 
% References and, for these, the correct style to use is ``Heading 5''. Use 
% ``figure caption'' for your Figure captions, and ``table head'' for your 
% table title. Run-in heads, such as ``Abstract'', will require you to apply a 
% style (in this case, italic) in addition to the style provided by the drop 
% down menu to differentiate the head from the text.

% Text heads organize the topics on a relational, hierarchical basis. For 
% example, the paper title is the primary text head because all subsequent 
% material relates and elaborates on this one topic. If there are two or more 
% sub-topics, the next level head (uppercase Roman numerals) should be used 
% and, conversely, if there are not at least two sub-topics, then no subheads 
% should be introduced.
